\section{helicity and transversity distributions}
In previous sections, we have considered the matching condition for unpolarized quark distribution. Here we present the results for the quark helicity and transversity distribution. For the quark helicity distribution in a longitudinally polarized quark, the quasi distribution $\Delta \tilde q^{(1)}(x)$ can be obtained by replacing $\gamma^z$ with $\gamma^z\gamma^5$ in Eq. (1). The one-loop result then reads
\begin{align}\label{helivertexnoexp}
\Delta \tilde q^{(1)}(x)&=\frac{\alpha_S C_F}{2\pi} \left\{ \begin{array} {ll} \frac{1+x^2}{1-x}\ln \frac{x(\Lambda(x)-x P^z)}{(x-1)(\Lambda(1-x)+P^z(1-x))}+1-\frac{x P^z}{\Lambda(x)}+\frac{x\Lambda(1-x)+(1-x)\Lambda(x)}{(1-x)^2 P^z}\ , & x>1\ , \\
\ \\
\frac{1+x^2}{1-x}\ln\frac{(P^z)^2}{m^2}+\frac{1+x^2}{1-x}\ln \frac{4x(\Lambda(x)-x P^z)}{(1-x)(\Lambda(1-x)+(1-x)P^z)}-\frac{4}{1-x}+2x+3\ & \\
-\frac{x P^z}{\Lambda(x)}+\frac{x\Lambda(1-x)+(1-x)\Lambda(x)}{(1-x)^2 P^z}\ , & 0<x<1\ . \\
\ \\
\frac{1+x^2}{1-x}\ln \frac{(x-1)(\Lambda(x)-x P^z)}{x(\Lambda(1-x)+(1-x)P^z)}-1-\frac{x P^z}{\Lambda(x)}+\frac{x\Lambda(1-x)+(1-x)\Lambda(x)}{(1-x)^2P^z}\ , & x<0\ . \end{array} \right.
\end{align}
Taking the limit $\mu\rightarrow \infty$ yields
\begin{align}\label{helivertexLgmu}
\Delta \tilde q^{(1)}(x)&=\frac{\alpha_S C_F}{2\pi} \left\{ \begin{array} {ll} \frac{1+x^2}{1-x}\ln \frac{x}{x-1}+1+\frac{\mu}{(1-x)^2P^z}\ , & x>1\ , \\ \frac{1+x^2}{1-x}\ln \frac{(P^z)^2}{m^2}+\frac{1+x^2}{1-x}\ln\frac{4x}{1-x}-\frac{4}{1-x}+2x+3+\frac{\mu}{(1-x)^2P^z}\ , & 0<x<1\ , \\ \frac{1+x^2}{1-x}\ln \frac{x-1}{x}-1+\frac{\mu}{(1-x)^2P^z}\ , & x<0\ . \end{array} \right.
\end{align}
The result for the light-cone distribution is again given by taking $P^z\to\infty$
\begin{align}\label{helivertexIMF}
\Delta \tilde q^{(1)}(x)&=\frac{\alpha_S C_F}{2\pi} \left\{ \begin{array} {ll} 0\ , & x>1\  \mbox{or}\  x<0\ , \\ \frac{1+x^2}{1-x}\ln \frac{\mu^2}{m^2}-\frac{1+x^2}{1-x}\ln{(1-x)^2}-\frac{2}{1-x}+2x\ , & 0<x<1\ .\end{array} \right.
\end{align}
Note that as in the unpolarized case, the collinear singularity in the quasi quark helicity distribution is exactly the same as in the light-cone distribution.


Similarly, for the transversity distribution in a transversely polarized quark, the quasi distribution $\delta \tilde q^{(1)}(x)$ is obtained by replacing $\gamma^z$ with $\gamma^z\gamma^\perp\gamma^5$ in Eq. (1). The one-loop result is
\begin{align}\label{transvertexnoexp}
\delta \tilde q^{(1)}(x)&=\frac{\alpha_S C_F}{2\pi} \left\{ \begin{array} {ll} \frac{2x}{1-x}\ln \frac{x(\Lambda(x)-xP^z)}{(x-1)(\Lambda(1-x)+P^z(1-x))}+\frac{x\Lambda(1-x)+(1-x)\Lambda(x)}{(1-x)^2 P^z}\ , & x>1\ , \\
\ \\
\frac{2x}{1-x}\ln\frac{(P^z)^2}{m^2}+\frac{2x}{1-x}\ln \frac{4x(\Lambda(x)-x P^z)}{(1-x)(\Lambda(1-x)+(1-x)P^z)}-\frac{4x}{1-x}+\frac{x\Lambda(1-x)+(1-x)\Lambda(x)}{(1-x)^2 P^z}\ , & 0<x<1\ , \\
\ \\
\frac{2x}{1-x}\ln \frac{(x-1)(\Lambda(x)-x P^z)}{x(\Lambda(1-x)+(1-x)P^z)}+\frac{x\Lambda(1-x)+(1-x)\Lambda(x)}{(1-x)^2P^z}\ , & x<0\ . \end{array} \right.
\end{align}
The limit $\mu\rightarrow \infty$ gives
\begin{align}\label{transvertexLgmu}
\delta \tilde q^{(1)}(x)&=\frac{\alpha_S C_F}{2\pi} \left\{ \begin{array} {ll} \frac{2x}{1-x}\ln \frac{x}{x-1}+\frac{\mu}{(1-x)^2P^z}\ , & x>1\ , \\
\frac{2x}{1-x}\ln \frac{(P^z)^2}{m^2}+\frac{2x}{1-x}\ln\frac{4x}{1-x}-\frac{4x}{1-x}+\frac{\mu}{(1-x)^2P^z}\ , & 0<x<1\ ,\\
\frac{2x}{1-x}\ln \frac{x-1}{x}+\frac{\mu}{(1-x)^2P^z}\ , & x<0\ , \end{array} \right.
\end{align}
and the result in the infinite momentum frame is
\begin{align}\label{transvertexIMF}
\delta q^{(1)}(x)&=\frac{\alpha_S C_F}{2\pi} \left\{ \begin{array} {ll} 0\ , & x>1\  \mbox{or}\  x<0\ , \\ \frac{2x}{1-x}\ln \frac{\mu^2}{m^2}-\frac{2x}{1-x}\ln{(1-x)^2}-\frac{2x}{1-x}\ , & 0<x<1\ .\end{array} \right.
\end{align}

One can construct similar matching conditions as in Eq. (12) for the helicity and transversity distributions.
We just list the results for the matching factors here, noting that the quark self-energy is the same.
For the quark helicity distribution, one has for $\xi>1$
\begin{equation}
   \Delta Z^{(1)}(\xi)/C_F = \left(\frac{1+\xi^2}{1-\xi}\right)\ln \frac{\xi}{\xi-1} + 1 + \frac{1}{(1-\xi)^2}\frac{\mu}{P^z} \ ,
\end{equation}
whereas for $0<\xi<1$
\begin{equation}
   \Delta Z^{(1)}(\xi)/C_F = \left(\frac{1+\xi^2}{1-\xi}\right)\ln \frac{(P^z)^2}{\mu^2}+\left(\frac{1+\xi^2}{1-\xi}\right)\ln\big[{4\xi(1-\xi)}\big]-\frac{2}{1-\xi}+3 +\frac{\mu}{(1-\xi)^2P^z} \ ,
\end{equation}
and for $\xi<0$
\begin{equation}
   \Delta Z^{(1)}(\xi)/C_F = \left(\frac{1+\xi^2}{1-\xi}\right)\ln \frac{\xi-1}{\xi}-1 +\frac{\mu}{(1-\xi)^2P^z} \ .
\end{equation}
The linearly divergent term is the same as in the unpolarized case.

Finally, in the factorization theorem for transversity distribution, one has the matching factor for $\xi>1$,
\begin{equation}
   \delta Z^{(1)}(\xi)/C_F = \left(\frac{2\xi}{1-\xi}\right)\ln \frac{\xi}{\xi-1} + \frac{1}{(1-\xi)^2}\frac{\mu}{P^z} \ ,
\end{equation}
whereas for $0<\xi<1$
\begin{equation}
   \delta Z^{(1)}(\xi)/C_F = \left(\frac{2\xi}{1-\xi}\right)\ln \frac{(P^z)^2}{\mu^2}+\left(\frac{2\xi}{1-\xi}\right)\ln\big[{4\xi(1-\xi)}\big]-\frac{2\xi}{1-\xi} +\frac{\mu}{(1-\xi)^2P^z} \ ,
\end{equation}
and for $\xi<0$
\begin{equation}
   \delta Z^{(1)}(\xi)/C_F = \left(\frac{2\xi}{1-\xi}\right)\ln \frac{\xi-1}{\xi}+\frac{\mu}{(1-\xi)^2P^z} \ .
\end{equation}
One again has an linearly divergent contribution. Near $\xi=1$, one needs to include an extra contribution from self energy as before.
